\documentclass[11pt,letterpaper]{article}

% ---------- Typography: serif body, sans headings (institutional memo standard) ----------
\usepackage[margin=1.05in,top=0.95in,bottom=1in]{geometry}
\usepackage[T1]{fontenc}
\usepackage{mathpazo}            % Palatino body — the classic finance-document serif
\usepackage{helvet}              % Helvetica for headings and labels only
\usepackage{microtype}           % refined kerning and justification
\usepackage{xcolor}
\usepackage[colorlinks=true,urlcolor=navy,linkcolor=navy]{hyperref}
\urlstyle{sf}                    % URLs in the document's sans font, not monospace
\usepackage{enumitem}
\usepackage{fancyhdr}
\usepackage{lastpage}

\definecolor{navy}{HTML}{1F3864}
\definecolor{midgray}{HTML}{595959}
\definecolor{hairline}{HTML}{C8C8C8}

\linespread{1.10}                % comfortable executive leading
\setlength{\parindent}{0pt}
\setlength{\parskip}{3.5pt}

% ---------- Footer: understated, informational ----------
\pagestyle{fancy}
\fancyhf{}
\renewcommand{\headrulewidth}{0pt}
\fancyfoot[L]{\sffamily\fontsize{7.5}{9}\selectfont\color{midgray}ENNACIRI \textbar\ EARLY-WARNING MODEL FOR U.S. BANK CAPITAL DISTRESS \textbar\ JULY 2026}
\fancyfoot[R]{\sffamily\fontsize{7.5}{9}\selectfont\color{midgray}\thepage\ OF \pageref{LastPage}}

% ---------- Section label: small uppercase bold sans, navy — corporate memo register ----------
\newcommand{\sectionhead}[1]{%
  \vspace{6pt}%
  {\sffamily\bfseries\fontsize{9.5}{12}\selectfont\color{navy}\MakeUppercase{#1}}\par
  \vspace{-9pt}\textcolor{hairline}{\rule{\textwidth}{0.5pt}}\par\vspace{0pt}}

% Sub-label within a section
\newcommand{\sublabel}[1]{{\sffamily\bfseries\fontsize{9}{11}\selectfont\color{midgray}#1}\par\vspace{2pt}}

% Answer body
\newcommand{\answer}[1]{#1\par\vspace{3pt}}

\begin{document}

% ---------- Title block ----------
\vspace*{-14pt}
{\sffamily\fontsize{8.5}{11}\selectfont\color{midgray}MSDS 696 DATA SCIENCE PRACTICUM II \quad\textbar\quad WEEK 1 PROJECT PROPOSAL}\par\vspace{10pt}
{\sffamily\bfseries\fontsize{21}{25}\selectfont\color{navy}Early-Warning Model for\\U.S. Bank Capital Distress}\par\vspace{6pt}
{\fontsize{12}{16}\selectfont\itshape\color{midgray}Identifying banks at risk of becoming undercapitalized while corrective actions are possible}\par\vspace{9pt}
{\sffamily\fontsize{9}{12}\selectfont Oussama Ennaciri \quad\textbar\quad July 6, 2026}\par\vspace{5pt}
\textcolor{navy}{\rule{\textwidth}{1.2pt}}

% ---------- Executive summary ----------
\sectionhead{Executive summary}
\answer{The 2023 failures of Silicon Valley Bank and First Republic could have been prevented if warning signs had been detected and acted on earlier. This project builds an early-warning model that identifies, from public quarterly regulatory data, which U.S.\ banks are at risk of becoming undercapitalized and how soon it could happen, while the bank can still take corrective action.}

% ---------- Required proposal fields ----------
\sectionhead{The problem, in one sentence}
\answer{U.S.\ banks that slide into undercapitalization lose their corrective options quarter by quarter, but the risk is visible in public regulatory data well before mandatory intervention; this project identifies which banks are at risk of becoming undercapitalized and how soon, while management can still act.}

\sectionhead{The three blanks a real problem fills}

\sublabel{The question}
\answer{Which U.S.\ banks are at risk of becoming undercapitalized, and how soon?}

\sublabel{The stakeholder: who acts on the answer}
\answer{\textbf{Decision-makers:}
\begin{itemize}[topsep=1pt,itemsep=1pt,leftmargin=14pt,label=\textcolor{navy}{\footnotesize\textbullet}]
  \item \textbf{The bank:} risk department, executive management, and the board.
  \item \textbf{The regulators:} FDIC, OCC, Federal Reserve, and state examiners.
\end{itemize}

\textbf{Affected parties:}
\begin{itemize}[topsep=1pt,itemsep=1pt,leftmargin=14pt,label=\textcolor{navy}{\footnotesize\textbullet}]
  \item \textbf{Inside the bank:} employees, shareholders, and the parent holding company, which must guarantee any capital restoration plan.
  \item \textbf{Direct counterparties:} depositors, especially uninsured; borrowers; creditors; the Federal Home Loan Banks; and deposit brokers.
  \item \textbf{The banking system:} peer banks exposed to contagion and higher FDIC assessments, the FDIC insurance fund, and acquiring institutions.
  \pagebreak
  \item \textbf{The public:} the Federal Reserve as lender of last resort, the U.S.\ Treasury, state and local governments, and, in systemic cases, taxpayers.
\end{itemize}}

\sublabel{The action: what they do differently once they have it}
\answer{\textbf{The bank:} raise capital, cut dividends, restructure the securities portfolio and deposit mix, shrink risky exposures, or pursue a voluntary merger.

\textbf{The regulators:} escalate examination frequency and apply supervisory tools before the tier drop; as a last resort, impose a capital restoration plan, dividend restrictions, and brokered-deposit prohibitions once options have narrowed; and when nothing else works, close the bank and place it in FDIC receivership, as happened to Silicon Valley Bank in 2023.}

\sectionhead{Personal angle}
\answer{Two years as a Financial Associate at a broker-dealer financial institution. I worked closely with the risk and compliance teams, helped prepare reports for NFA audits, and saw firsthand what the auditors flagged and why. I also passed the SIE exam, so I know the regulatory basics behind capital requirements, and I can defend which balance-sheet ratios matter and what a risk committee actually does when a warning shows up.}

\sectionhead{Candidate data}
\answer{\begin{itemize}[topsep=1pt,itemsep=1pt,leftmargin=14pt,label=\textcolor{navy}{\footnotesize\textbullet}]
  \item \textbf{FDIC BankFind Suite API} (primary source): quarterly financial data for every FDIC-insured institution, 1992 to present, including the Prompt Corrective Action capital ratios used to define the target.
  \item \textbf{FRED} (supporting source): macroeconomic context, such as interest rates, unemployment, and other indicators as the analysis develops.
\end{itemize}}

\sectionhead{First milestones / rough plan}
\answer{\begin{enumerate}[topsep=1pt,itemsep=1pt,leftmargin=16pt,label=\textcolor{navy}{\bfseries\arabic*.}]
  \item Literature review of bank failure and existing early-warning models (\textbf{in progress})
  \item Collect the data (FDIC API, 1992 to present)
  \item Understand the data through exploratory data analysis
  \item Prepare the data for machine learning
  \item Train and compare three candidate models
\end{enumerate}}

\begin{samepage}
\sectionhead{Track}
\answer{Own idea. Data verified as real and reachable (live FDIC API queries, July 2026). Fallback: curated project \#7, CFPB consumer-complaint triage.}
\end{samepage}

\vspace{2pt}
{\small\itshape\color{midgray}\textbf{Note on LLM use:} This document was typeset in LaTeX using Claude Code, which handled the formatting, styling code, and wording polish. All content, research, and decisions are my own, and every suggested edit was reviewed and approved by me.}

\end{document}
